% octavo — a modification of the standard `book` class implementing
% Tschichold/Rosarivo classical page-proportion and margin theory, plus
% signature imposition for hand bookbinding. Ships with TeX Live; not
% bundled.
%
% `demy` selects one of the class's nine classical paper formats
% (foolscap/crown/post/largepost/demy/medium/royal/superroyal/imperial).
% `titlepage` is passed explicitly: the class's own \ExecuteOptions sets
% notitlepage as the ACTUAL default (a standalone \@titlepagetrue earlier
% in the file is overridden by it), so \maketitle is undefined unless
% this is passed.
%
% Class documentation: texdoc octavo | https://ctan.org/pkg/octavo
% Licence: LPPL, (c) 2003, 2007 Stefan A. Revets — the grant is in the
% package's octavo.ins file, not octavo.dtx itself.
\documentclass[demy,11pt,openright,titlepage]{octavo}

\usepackage[utf8]{inputenc}
\usepackage[T1]{fontenc}

\title{On the Proportions of the Well-Made Book}
\author{Cecilia Hartmann}
\date{\today}

\begin{document}

\frontmatter

\maketitle

\chapter*{Preface}
This short treatise gathers a printer's working notes on classical
page design — the kind of book Jan Tschichold or Hugh Williamson might
have kept beside the composing stone. It is set in the \emph{demy}
format, one of nine classical trim sizes this class provides.

\tableofcontents

\mainmatter

\chapter{The Golden Canon}

Long before typography had a name, scribes and printers noticed that
certain page proportions felt more settled than others. The
\emph{golden canon of page construction}, reconstructed by Jan
Tschichold from medieval manuscripts, divides the page and its text
block so that the margins grow in a fixed ratio outward from the
spine: the inner margin narrowest, the head somewhat wider, the outer
wider still, and the foot widest of all.

\section{Why the Margins Are Not Equal}

An equal margin on all four sides looks correct on a single sheet, but
wrong in a bound book: two facing pages with equal margins read as a
single wide gutter, pulling the eye toward the spine instead of toward
the text. Tschichold's unequal margins compensate for the way an open
book is actually seen — as a spread, not as two isolated rectangles.

\chapter{Signatures and Imposition}

A book is not printed one page at a time. Pages are grouped into
\emph{signatures} — folded sheets of 8, 16, or 32 pages — and the
pages within a signature are arranged, or \emph{imposed}, so that
after folding and trimming they fall into the correct reading order.

\section{Folding an Octavo Signature}

An octavo signature folds a single large sheet in half three times,
producing eight leaves (sixteen pages) per signature — hence the class
and format's shared name. Getting the imposition wrong is invisible on
screen and catastrophic on paper: pages print in the wrong order, or
upside down, only discovered after folding.

\backmatter

\chapter{Further Reading}

Tschichold's \emph{The Form of the Book} and Williamson's \emph{Methods
of Book Design} remain the standard references for the ideas this
class implements in code.

\end{document}
